% Conclusions section

We have developed a data-efficient machine-learning framework for predicting the formation energies and sublattice occupancies of complex TCP intermetallic phases in the Fe--Mo binary system. By encoding domain knowledge of chemistry, crystallography, and local bonding into the feature representation, we achieve accurate predictions using only 262 DFT training structures of simple TCP phases.

Key findings include:

\begin{enumerate}
  \item BOP descriptors, derived from a canonical tight-binding model, provide the most transferable and accurate features for TCP energy prediction, outperforming ACE and SOAP across all test and validation sets.
  \item Coordination-number-resolved averaging (CNAV) improves prediction accuracy by 20--40\% compared with conventional structure-level averaging, demonstrating the value of explicit crystallographic knowledge.
  \item The VotingRegressor ensemble of KRR, RF, and MLP models ensures robust predictions, with ensemble averaging providing greater benefit in the small-data regime.
  \item Predicted formation energies, when used within the Bragg--Williams--CEF thermodynamic framework, yield sublattice occupancies consistent with experimental X-ray diffraction data for the R phase.
\end{enumerate}

This work establishes a pathway for efficient computational discovery of TCP phase stability in multi-principal element alloys, where the combinatorial explosion of possible configurations makes exhaustive DFT sampling infeasible.
